I have everything I need. Here is the LaTeX body for the assigned subsection.

\subsection{Machine Learning Pipeline}

Every 0.40\,s window (200 samples at 500\,Hz, 60 channels) must yield three answers at once: is any leg entangled, which leg, and how severe. A model built for this has to respect one hard operational constraint. At inference the detector sees a stream that ends at the current instant and can never peek ahead, so the network must produce, at time $t$, a decision that depends only on samples up to and including $t$. We use a temporal convolutional network (TCN) \cite{bai2018tcn} because it satisfies this constraint by construction, processes a whole window in parallel rather than one step at a time, and grows its temporal reach cheaply through dilation. The architecture is shown in Fig.~\ref{fig:tcn}, and its hyperparameters are collected in Table~\ref{tab:hparams}.

Causality is enforced through left-padding. A standard convolution centred on $t$ mixes in future samples; instead we pad only the left (past) side of each sequence so that the output at $t$ is a function of $\{x_{t-k},\dots,x_t\}$ and nothing later. The immediate benefit is that the model we train offline on complete recorded windows is numerically the same function that runs online on a growing buffer: there is no train-test discrepancy from windowing or look-ahead, a property we verify empirically later. This mirrors the causal-convolution idea from WaveNet \cite{oord2016wavenet}, adapted here to multivariate proprioception rather than audio.

The input first passes through a causal convolutional stem that lifts the 60 input channels to a 64-dimensional feature space, giving the residual stack a consistent width to operate on. Five residual blocks follow, with dilation factors $1,2,4,8,16$. Dilation inserts gaps between the taps of a kernel so that a fixed kernel size covers an exponentially longer span as we stack blocks, which is how a small network reaches far back in time without a matching blow-up in parameters or compute. With kernel size 3 and these five dilations, the receptive field is 127 samples, or 254\,ms at 500\,Hz. That is comfortably inside the 0.40\,s window, so the last output position has actually seen most of a trot stride before it commits to a decision.

Each residual block has the same internals: a causal convolution, a GELU nonlinearity \cite{hendrycks2016gelu}, a second causal convolution, another GELU, dropout at rate 0.1, and then an identity skip connection that adds the block input back to its output. The skip connection lets gradients flow past the block and makes each block learn a refinement of its input rather than a full remapping, which stabilises training of the stack. Both convolutions use weight normalization \cite{salimans2016weightnorm}, which decouples the scale of each filter from its direction and tends to smooth optimization. GELU is a smooth gating nonlinearity that we found preferable to a hard rectifier for these continuous torque and velocity signals.

After the stack we read out only the embedding at the last timestep, a single 64-dimensional vector. This is deliberate: at deployment the quantity of interest is the decision \emph{now}, given the history in the buffer, and the last position is precisely the one whose receptive field ends at the present sample. That 64-D readout feeds three linear heads. The detection head emits one logit (entangled or not). The leg head emits four logits, treated as independent labels because more than one leg can be affected. The intensity head emits four values, one per leg. The whole network has 136{,}393 parameters, roughly 136k, small enough to ship as a half-megabyte model.

The three tasks are trained jointly with a single weighted objective. Detection and leg attribution are classification problems and use binary cross-entropy; intensity regression uses a Huber loss, which behaves like squared error for small residuals but like absolute error for large ones and so is less swayed by outliers. Detection and leg are weighted equally, and the intensity term is down-weighted because it is an auxiliary signal, not a headline output. In words, the loss is one part detection cross-entropy, one part leg cross-entropy, and three-tenths of a part intensity Huber:
\begin{equation}
\mathcal{L} = 1.0\,\mathrm{BCE}_{\text{det}} + 1.0\,\mathrm{BCE}_{\text{leg}} + 0.3\,\mathrm{Huber}_{\text{int}} .
\end{equation}
The intensity term deserves a caveat. Its Huber ($\delta=0.1$) is computed against a physics pseudo-target, not a human-labeled severity, and elements corresponding to the affected leg are up-weighted $3{:}1$ relative to the rest. This is an up-weight, not a hard mask: because most legs at most times carry a target of zero, leaving the weighting flat would let ``predict zero everywhere'' dominate the loss and the head would learn nothing about the affected leg. The head therefore ends up as a weak learner that tracks the physics target, while the physics formula itself remains authoritative at inference.

Two further measures handle the strong class imbalance between normal and entangled windows. A class-balanced weighted sampler equalises exposure to positives and negatives during training, and each classification task carries a (capped) positive-class weight in its cross-entropy. Optimization uses AdamW \cite{loshchilov2019adamw} with a cosine-annealed learning rate over 30 epochs; the exact learning rate, weight decay, batch size, and seed are listed in Table~\ref{tab:hparams}.

Severity is reported through a physics-grounded index rather than the network head alone, because no ground-truth severity labels exist to regress against. The index combines three quantities per leg. A confidence gate $g$ is the per-leg entanglement probability, so the index is near zero whenever the leg is not flagged and cannot manufacture severity out of noise. A deviation term $d$ is the Mahalanobis distance of the window's per-leg feature vector from the normal walking distribution (mean and covariance fit on training walking data only), measuring how far the leg's kinematics have drifted from ordinary gait. A resistive-effort term $r$ is the mean of the downward thigh torque divided by thigh joint speed, high when a leg pushes hard yet barely moves, the signature of a snag. These combine as
\begin{equation}
I_{\text{leg}} = g\,\sqrt{d\,r},
\end{equation}
with the square root compressing the product of deviation and effort, and the result calibrated to $[0,1]$. Because there is no severity ground truth, we present $I_{\text{leg}}$ as a calibrated index, not a measured force. The value deployed to the runtime blends the network head and this physics index in equal parts, so learned evidence and analytic evidence are given the same say.

The choice of a TCN over the obvious alternatives follows from these requirements. A recurrent network is sequential and awkward to make strictly causal-equivalent between offline and online use, and it trains one step at a time; the TCN is causal by padding, evaluates a window in parallel, and reaches 254\,ms of history for a small parameter budget. A gradient-boosted tree ensemble, though a strong tabular baseline, has no native notion of time and would have to be handed hand-built window summaries, discarding exactly the temporal structure that distinguishes a transient snag from a steady stance.